\kchapter{Organising: Labels, Selectors, Annotations, Namespaces}{ns}

\begin{objectives}
\begin{itemize}
  \item how labels and selectors glue Kubernetes objects together;
  \item the recommended label set;
  \item when to use annotations instead of labels;
  \item how namespaces divide a cluster, and what they do \emph{not} isolate.
\end{itemize}
\end{objectives}

\section{Labels: sticky notes on objects}

A \kterm{label} is a key/value pair attached to an object's \kc{metadata}.
Labels mean nothing to Kubernetes on their own --- \emph{you} decide what they
mean --- but they are how objects find each other.

\begin{yamlcode}{labels.yaml}
metadata:
  name: api-7d9f8c-xk2lp
  labels:
    app.kubernetes.io/name: taskflow-api
    app.kubernetes.io/part-of: taskflow
    app.kubernetes.io/version: "1.4.0"
    app.kubernetes.io/component: backend
    environment: production
    team: payments
\end{yamlcode}

\section{Selectors: queries over labels}

A \kterm{selector} picks objects by their labels. You will see selectors
everywhere:

\begin{itemize}
  \item a \textbf{Service} sends traffic to Pods whose labels match its selector;
  \item a \textbf{Deployment/ReplicaSet} owns the Pods matching its selector;
  \item a \textbf{NetworkPolicy} applies to Pods matching its selector;
  \item \kc{kubectl} filters with \kc{-l}.
\end{itemize}

\begin{bashcode}{terminal}
kubectl get pods -l app=web                       # equality
kubectl get pods -l 'environment in (staging,production)'   # set-based
kubectl get pods -l 'app=web,tier!=cache'          # AND, not-equal
kubectl get pods --show-labels
kubectl get pods -L app,version                    # labels as columns
\end{bashcode}

In YAML, selectors come in two flavours:

\begin{yamlcode}{selectors.yaml}
# Simple form (Services)
selector:
  app: web

# Rich form (Deployments, Jobs, NetworkPolicies...)
selector:
  matchLabels:
    app: web
  matchExpressions:
    - key: tier
      operator: In            # In, NotIn, Exists, DoesNotExist
      values: [frontend, edge]
\end{yamlcode}

\begin{center}
\begin{tikzpicture}[
  podn/.style={fill=k8sblue!12, draw=k8sblue, rounded corners=4pt, align=center,
               font=\scriptsize, minimum width=2.4cm, minimum height=1.1cm},
  podx/.style={fill=k8sgray!10, draw=k8sgray, rounded corners=4pt, align=center,
               font=\scriptsize, minimum width=2.4cm, minimum height=1.1cm, text=k8sgray},
  >={Stealth[length=2.6mm]}]
  \node[fill=k8sblue, text=white, rounded corners=5pt, align=center, font=\small\bfseries,
        minimum width=3cm, minimum height=1.1cm] (svc) at (0,2.4)
        {\kicon{svc}\ Service\\[-1pt]\scriptsize\mdseries selector: app=web};
  \node[podn] (p1) at (-3.6,0) {\textbf{Pod}\\app=web\\tier=frontend};
  \node[podn] (p2) at (-0.9,0) {\textbf{Pod}\\app=web\\tier=frontend};
  \node[podn] (p3) at (1.8,0)  {\textbf{Pod}\\app=web\\tier=edge};
  \node[podx] (p4) at (4.5,0)  {\textbf{Pod}\\app=db};
  \foreach \p in {p1,p2,p3} \draw[->, thick, k8sblue] (svc) -- (\p);
  \draw[->, thick, k8sgray, dashed] (svc) -- node[right, font=\scriptsize]{\faTimes} (p4);
\end{tikzpicture}
\end{center}

\begin{warningbox}
A Deployment's \kc{spec.selector} is \textbf{immutable} and must match its Pod
template labels. Choose stable labels (like \kc{app.kubernetes.io/name}), never
something that changes per release like a version number.
\end{warningbox}

\section{Annotations: notes for tools and humans}

\kterm{Annotations} are also key/value metadata, but they are \emph{not} used
for selecting. Values can be long (even JSON). Tools read them to change
behaviour.

\begin{yamlcode}{annotations.yaml}
metadata:
  annotations:
    kubernetes.io/change-cause: "upgrade to 1.4.0 for faster search"
    prometheus.io/scrape: "true"
    owner-contact: "payments-oncall@example.com"
    description: "Public API for task management"
\end{yamlcode}

\begin{center}
\ktablestart
\begin{tabularx}{\linewidth}{P{4.2cm} L L}
\kheader \kth{} & \kth{Labels} & \kth{Annotations} \\
Used by selectors & Yes & No \\
Value size & Short (63 chars) & Large (up to 256\,KB total) \\
Typical use & Identity, grouping, routing & Tool configuration, docs, audit trail \\
\end{tabularx}
\end{center}

\section{Namespaces: rooms in the house}

A \kterm{namespace} is a virtual partition of a cluster. Names need only be
unique \emph{within} a namespace, so team A and team B can both have a
Deployment called \kc{api}.

\begin{bashcode}{terminal}
kubectl get namespaces
# default, kube-system, kube-public, kube-node-lease  (+ your own)

kubectl create namespace dev
kubectl apply -f app.yaml -n dev
kubectl get pods -n dev
kubectl get pods --all-namespaces      # or -A
kubectl config set-context --current --namespace=dev
\end{bashcode}

\begin{yamlcode}{namespace.yaml}
apiVersion: v1
kind: Namespace
metadata:
  name: dev
  labels:
    team: platform
    pod-security.kubernetes.io/enforce: baseline   # see Chapter 16
\end{yamlcode}

\subsection{What namespaces give you}

\begin{itemize}
  \item a scope for names;
  \item a boundary for \textbf{RBAC} permissions (Chapter 16);
  \item a boundary for \textbf{ResourceQuotas} and \textbf{LimitRanges} (Chapter 14);
  \item a target for \textbf{NetworkPolicies} (Chapter 11);
  \item a DNS segment: \kc{api.dev.svc.cluster.local}.
\end{itemize}

\begin{warningbox}
By default, namespaces are \textbf{not} a network or security wall. A Pod in
\kc{dev} can talk to a Pod in \kc{prod} unless you add NetworkPolicies, and
Pods on the same node share a kernel. For hard multi-tenancy, use separate
clusters or strong policy tooling.
\end{warningbox}

\subsection{Cluster-scoped objects}

Some objects do not live in any namespace: Nodes, Namespaces themselves,
PersistentVolumes, StorageClasses, ClusterRoles, CRDs. Check with
\kc{kubectl api-resources --namespaced=false}.

\subsection{A common layout}

\begin{itemize}
  \item one namespace \emph{per application per environment} (\kc{shop-dev},
        \kc{shop-staging}) on shared clusters, or
  \item one namespace \emph{per application} with separate clusters per
        environment (recommended for production isolation);
  \item platform tools in their own namespaces: \kc{monitoring},
        \kc{argocd}, \kc{cert-manager}, \kc{ingress}.
\end{itemize}

\begin{recap}
\begin{itemize}
  \item Labels identify and group; selectors query labels.
  \item Services, Deployments, and policies all find Pods via selectors.
  \item Annotations carry non-identifying metadata for tools.
  \item Namespaces scope names, RBAC, and quotas --- but are not a firewall by themselves.
\end{itemize}
\end{recap}
